\begin{tikzpicture}[x=1mm,y=1mm,font=\figurefont\footnotesize,
  every node/.style={text=NNInk},
  nn flow/.append style={draw=NNWire,rounded corners=1.5mm},
  nn operator path/.append style={draw=NNWire,rounded corners=1.5mm},
  op/.style={nn operator,minimum width=43mm,minimum height=10mm,inner sep=2pt},
  token/.style={draw=NNInput,fill=NNInput!18,line width=.85pt,minimum width=6mm,minimum height=7mm,inner sep=1pt}]
  \node[nn heading] at (23,17) {卷积前端：缩短序列};
  \node[nn heading] at (82,17) {交替：状态与显式读取};
  \node[nn heading] at (141,17) {并行：融合两类表示};
  \foreach \i in {1,...,6}{\node[token] (a\i) at ({8*\i-5},0) {$\i$};}
  \foreach \i in {1,2,3}{
    \pgfmathtruncatemacro{\lefttoken}{2*\i-1}
    \pgfmathtruncatemacro{\righttoken}{2*\i}
    \node[token,draw=NNHidden,fill=NNHidden!18,minimum width=12mm] (c\i) at ({16*\i-9},-24) {$\lefttoken{:}\righttoken$};
    \draw[nn operator path] (a\lefttoken.south)--(c\i.north west);
    \draw[nn operator path] (a\righttoken.south)--(c\i.north east);
  }
  \node[nn annotation] at (23,9) {窗口2，步长2};
  \node[op] (att0) at (23,-48) {注意力块};
  \foreach \i in {1,2,3}{\draw[nn operator path] (c\i.south)--({16*\i-9},-43);}
  \node[nn annotation,align=center] at (23,-69) {6个输入位置\\变成3个特征位置};
  \foreach \i/\word in {1/猫,2/追,3/狗}{
    \node[token,minimum width=11mm] (b\i) at ({16*\i+50},0) {\word};
    \node[op,minimum width=11mm] (s\i) at ({16*\i+50},-23) {SSM};
    \draw[nn flow] (b\i.south)--(s\i.north);
  }
  \draw[nn operator path] (s1.east)--node[above] {$\bm s_1$} (s2.west);
  \draw[nn operator path] (s2.east)--node[above] {$\bm s_2$} (s3.west);
  \node[op] (att1) at (82,-48) {注意力块};
  \foreach \i in {1,2,3}{\draw[nn operator path] (s\i.south)--({16*\i+50},-43);}
  \node[op] (ssm2) at (82,-73) {下一层SSM};
  \draw[nn operator path] (att1.south)--(ssm2.north);
  \node[token,minimum width=36mm] (seq) at (141,0) {猫\quad 追\quad 狗};
  \node[op,minimum width=21mm] (local) at (128,-24) {局部卷积};
  \node[op,minimum width=21mm] (global) at (154,-24) {注意力};
  \draw[nn flow,-,sharp corners] (seq.south)--(141,-10);
  \draw[nn flow] (141,-10)--(128,-10)--(local.north);
  \draw[nn flow] (141,-10)--(154,-10)--(global.north);
  \fill[NNWire] (141,-10) circle[radius=.45mm];
  \foreach \x/\name in {128/local,154/global}{
    \foreach \i in {0,1,2}{\draw[draw=NNHidden,fill=NNHidden!18,line width=.85pt] ({\x-9+6*\i},-39) rectangle ({\x-3+6*\i},-46);}
    \foreach \i/\word in {0/猫,1/追,2/狗}{\node at ({\x-6+6*\i},-42.5) {\word};}
    \draw[nn operator path] (\name.south)--(\x,-39);
  }
  \node[nn transform,minimum width=38mm] (merge) at (141,-61) {相加／拼接后投影};
  \draw[nn operator path] (128,-46)--(128,-51)--(merge.north west);
  \draw[nn operator path] (154,-46)--(154,-51)--(merge.north east);
  \node[nn annotation] at (141,-79) {仍与猫、追、狗逐位置对齐};
\end{tikzpicture}
